\section*{EMSx programs, revision of November 1, 2026}

This repository contains calculation programs, experiment configurations, tests,
and selected aggregate results for probabilistic correction of issued photovoltaic
forecasts. The current revision is \texttt{round1-2026-11-01}. The manuscript,
reviewer comments, response letter, raw measurements, full prediction arrays, and
trained model weights are excluded. The EMSx source data are available at
\url{https://doi.org/10.5281/zenodo.5510400}.

\section*{What changed}

The original WIS implementation paired interval bounds with reversed miscoverage
levels. All four scoring entry points now use one shared implementation in
\texttt{scripts/emsx\_probabilistic.py}. It uses the median and central intervals
covering 20, 40, 60, and 80 percent, with miscoverage levels 0.8, 0.6, 0.4, and 0.2,
respectively. The standard weights are described by Bracher et al.,
\url{https://doi.org/10.1371/journal.pcbi.1008618}. The outer interval covering
90 percent is evaluated separately. Tests independently check the equivalence
with the nine inner pinball losses.

Recoverable historical WIS values have been corrected. Values that cannot be
reconstructed from the retained aggregates are null in JSON or empty in CSV,
with their earlier numbers preserved under names beginning with
\texttt{legacy\_invalid\_}. Those earlier numbers are not valid WIS results.
\texttt{results/SCORING\_CORRECTIONS.json} records the changes. Complete original
run records containing unrounded pinball loss and the outer interval score can
recover WIS without retraining; the summary and verification programs apply
that identity. Aggregate means alone cannot recover missing seed variation.

The original partitions are chronological within each system but have overlap
across systems during global fitting. They therefore remain retrospective
comparisons. The new calendar experiment uses common cutoffs on January 1,
April 1, and July 1, 2016. Eligibility leaves 47 systems and 1,392,864 test
deliveries. The point comparison includes LightGBM, Transformer, BiLSTM, vendor,
seasonal naive, and persistence forecasts. All three neural seeds are retained.
LightGBM refits on training and tuning, while neural models retain training
checkpoints selected on tuning, so the budgets differ.

The quantile check reuses the common point tuning settings, the existing eleven
quantiles, and the window of 30 days. This check was added after inspection of
the point comparison, without selecting settings on quantile test outcomes.
Foundation adaptation has not been repeated under these common calendar cutoffs.
The original adaptation settings used original test outcomes and remain exploratory.

\section*{Environment and data}

Run commands from the repository directory, with inputs and newly computed
outputs in an external working directory. Python 3.12 was used. Install
\texttt{requirements-round1.txt} in a dedicated environment for the new experiments.
The original \texttt{requirements.txt} and the separate Chronos and Moirai
requirements remain available for the historical experiments. Their environments
should not be combined. Package versions and configuration values are recorded
in \texttt{results/round1/software\_versions.json} and
\texttt{protocols/emsx\_round1\_calendar\_20261101.json}.

The following shell commands show the data setup on Linux. Replace the example
working directory before running them.

\begin{verbatim}
export EMSX_WORK_ROOT=/absolute/path/to/emsx-work
python3 -m venv "$EMSX_WORK_ROOT/.venv-round1"
source "$EMSX_WORK_ROOT/.venv-round1/bin/activate"
python -m pip install -r requirements-round1.txt
python scripts/download_emsx.py
python scripts/build_emsx_daily_curves.py \
  --raw-root "$EMSX_WORK_ROOT/data/raw/emsx" \
  --output-root "$EMSX_WORK_ROOT/data/processed/emsx/daily_curves" \
  --manifest "$EMSX_WORK_ROOT/results/daily_profiles.json"
\end{verbatim}

On Windows, set \texttt{EMSX\_WORK\_ROOT} in PowerShell and invoke the Python
executable from the environment for every command. Activation is optional.

\begin{verbatim}
$env:EMSX_WORK_ROOT = 'C:\EMSx-work'
python -m venv "$env:EMSX_WORK_ROOT\.venv-round1"
$emsxPython = "$env:EMSX_WORK_ROOT\.venv-round1\Scripts\python.exe"
& $emsxPython -m pip install -r requirements-round1.txt
\end{verbatim}

By default, prepared profiles are read from the working directory under
\texttt{data/processed/emsx/daily\_curves}, and new outputs go under
\texttt{round1}. Optional environment variables \texttt{EMSX\_CURVE\_ROOT},
\texttt{EMSX\_PREDICTION\_ROOT}, and \texttt{EMSX\_ROUND1\_ROOT} override these
locations. None depends on the author's manuscript directory or server paths.

\section*{Repeat the calendar experiment}

Use the environment's Python executable in the following commands. LightGBM
must finish before the quantile experiment reads its selected settings, and all
six neural fits must finish before the combined summary is generated.

\begin{verbatim}
python scripts/calendar_validation.py --model lightgbm
python scripts/calendar_validation.py --model neural
python scripts/calendar_quantiles.py
python scripts/analyse_emsx_calendar.py
python scripts/verify_round1.py \
  --prediction-root "$EMSX_WORK_ROOT/round1/results" \
  --output "$EMSX_WORK_ROOT/round1/verification.json"
\end{verbatim}

Training and prediction can take substantial time on a CPU. The published
programs retain the calculation settings used for the saved results. Their
portable paths, output handling, and scoring code were checked against the
saved predictions; packaging did not retrain every model. Floating point and
training results can vary with platform and numerical libraries. The verifier
uses a tolerance of 0.001 kWh for point summaries and a stricter tolerance for
the quantile scores.

\section*{Check results without training}

The bundled summaries can be checked without downloading models or running
training. The first command below checks aggregate consistency, not full
predictions. Supplying the external prediction directory enables checks of
both partitions, every neural seed, matched delivery keys, finite values,
and recalculated scores.

\begin{verbatim}
python scripts/verify_round1.py
python -m unittest discover -s tests -v
python scripts/check_distribution.py
\end{verbatim}

The optional \texttt{EMSX\_TEST\_TMP} variable selects a writable directory for
temporary test files when the default temporary directory is unavailable.

\texttt{scripts/audit\_revision.py} reproduces the chronology audit and corrected
scores from the original prediction files. \texttt{scripts/check\_calibration.py}
checks the causal offsets at the first original test issue of each system.
These programs need the original prepared profiles and prediction files in the
external working directory. The original inference programs regenerate those
files when the EMSx data and required model checkpoints are available.

\section*{Included numerical results}

\texttt{results/round1} contains the common calendar experiment, numerical audit,
sample dates, software versions, and comparison figure. Vendor test MAE is
100.57 kWh, compared with 83.26 kWh for residual LightGBM, 89.35 kWh for BiLSTM,
and 90.66 kWh for the Transformer. The latter two are means across three seeds.
Rolling calibration changes WIS from 66.80 to 64.44 kWh and interval coverage
from 0.826 to 0.915, with wider intervals. Original sample results remain in the
other result directories, with the scoring corrections clearly recorded.

The repository does not establish transfer to other vendors or climate zones.
The calendar experiment reuses EMSx records examined during study development
and is not a prospective replication. The bootstrap resamples systems and daily
profiles rather than preserving dependence across successive days.

\section*{Second revision, November 2, 2026}

Version \texttt{round2-2026-11-02} adds the programs of the additional analyses
of the second revision. \texttt{scripts/round2\_compute.py} runs the leakage
comparison, the feature ablation, the ridge references, the cross-system error
feature, the transfer to systems excluded from fitting, and the block bootstrap.
\texttt{scripts/round2\_overlap\_check.py} and \texttt{scripts/round2\_drift\_check.py}
check the overlap of the partitions and the drift of the number of systems.
\texttt{scripts/round2\_restore\_analyses.py} recomputes the regime, season, and
reliability analyses on the common calendar test period. The selected LightGBM
settings and seeds are fixed in the programs. Aggregate numerical summaries are in
\texttt{results/round2}. The prediction arrays and the source data are not stored
here. The first revision remains identified by the tag \texttt{round1-2026-11-01}.
